\section{Architecture B --- Existing Vehicle / Alfred Gateway}
\label{sec:arch-b}

Customer already has an established vehicle E/E architecture and does not want to redesign the network, replace existing ECUs, replace CAN/LIN, or start a new vehicle program solely for Alfred. Alfred attaches itself to the existing infrastructure; the original network continues operating unmodified.

\begin{diagram}
                EXISTING VEHICLE

      ECU -------- CAN -------- ECU
                    |
                    |
              Alfred Node
              Gateway Mode
                    |
               10BASE-T1S
                    |
                    v
                 Clicker 4
\end{diagram}

\subsection{The Alfred pseudo-bus}

Alfred does not replace the original CAN/LIN bus. It builds a parallel, Alfred-facing representation of it:

\begin{diagram}
             EXISTING BUS
                 |
                 v
               CAPTURE
                 |
                 v
              TIMESTAMP
                 |
                 v
           NORMALIZE / MAP
                 |
                 v
               MIRROR
                 |
                 v
           10BASE-T1S
                 |
                 v
          ALFRED PSEUDO-BUS
\end{diagram}

Five functions, kept explicitly distinct rather than blurred together:

\begin{center}
\small
\begin{tabularx}{\linewidth}{@{}p{2.6cm} X@{}}
\toprule
\textbf{Function} & \textbf{Definition} \\
\midrule
Capture & Observing existing traffic. Read-only at the transceiver level; the only function active in Mode~1 (Section~\ref{sec:gateway-modes}). \\
Translation & Converting bus representation (CAN frame $\rightarrow$ normalized, timestamped record) -- a data-format operation, not a transmission. \\
Mirroring & Replicating useful, normalized traffic onto Alfred's own T1S network so the Clicker~4 (or any future central compute) sees legacy-bus state without touching the legacy bus itself. \\
Replay & Reproducing previously observed traffic during an authorized test, from a recording -- used for bench validation, not for live vehicle operation. \\
Active translation & Receiving an Alfred-side command and transmitting the corresponding \emph{validated} command onto the existing bus -- the only function that ever puts a new frame on the legacy bus, and the only one gated by Mode~2. \\
\bottomrule
\end{tabularx}
\end{center}

\subsection{Authorized bidirectional gateway}
\label{sec:gateway-modes}

The Gateway supports two explicit hardware operating states, not a software flag alone.

\begin{diagram}
MODE 1 -- PASSIVE (default at power-up)
        Vehicle ---------> Alfred
        Vehicle <----X---- Alfred        (TX physically gated off)

MODE 2 -- AUTHORIZED BIDIRECTIONAL
        Vehicle <--------> Alfred
   Alfred -> 10BASE-T1S -> Gateway -> translation -> CAN/CAN-FD/LIN -> vehicle
\end{diagram}

Mode~1 is the safe default; the node must not transmit onto the legacy bus, and where practical this is enforced with a hardware TX-enable gate (transceiver enable pin held off by default), not only by not calling a software transmit function. Mode~2 additionally includes: a physical TX-enable line, watchdog-supervised command timeout (a stale authorization reverts to Mode~1), bus-off detection and handling, error counters, a clear physical/visual indication of current mode, a separate debug enable, an append-only audit log of every frame the Gateway ever transmits, and deterministic reset behavior (reset always re-enters Mode~1). This capability is assumed to be used only on Alfred-owned hardware, customer development vehicles, test benches, or explicitly authorized customer systems -- circumventing OEM security controls is out of scope categorically, not a design constraint to work around.

\subsection{Brownfield bus support priority}

\begin{center}
\small
\begin{tabularx}{\linewidth}{@{}p{3.2cm} X@{}}
\toprule
\textbf{Tier} & \textbf{Buses} \\
\midrule
Mandatory for November & CAN, CAN-FD, 10BASE-T1S \\
Desired for November & LIN \\
Optional, after November & 100BASE-T1, 1000BASE-T1, specialized/proprietary buses \\
\bottomrule
\end{tabularx}
\end{center}

\begin{decision}[title={LIN in Rev~C: schedule-conditional, not guaranteed}]
Include a LIN transceiver footprint on Rev~B (Section~\ref{sec:rev-b}) at negligible incremental schedule cost, since it shares the same MCU UART resource strategy as CAN-FD bring-up. Whether LIN is \emph{validated and demonstrated} in the November~30 gate depends on Rev~B bring-up going to schedule; if CAN/CAN-FD/T1S validation is behind at the mid-November checkpoint (Section~\ref{sec:schedule}), LIN demonstration is cut first (Section~\ref{sec:feature-cuts}) -- CAN/CAN-FD/T1S validation is never sacrificed to add protocol count.
\end{decision}

\subsection{Brownfield November demonstration}

\begin{diagram}
       Device A
          |
         CAN
          |
       Device B
          |
          +-------- Alfred Gateway
                       |
                  10BASE-T1S
                       |
                    Clicker 4

PHASE A -- PASSIVE
  Physical event -> legacy bus traffic -> Gateway captures it ->
  traffic mirrored onto T1S -> Clicker 4 receives it

PHASE B -- BIDIRECTIONAL
  Clicker 4 -> T1S command -> Gateway -> legacy CAN command ->
  legacy device -> physical response -> observed and mirrored back
\end{diagram}

A controlled legacy network (two devices on a CAN segment, on the bench, not a real vehicle) is built specifically for this demonstration -- the point is validated gateway behavior, not sourcing an actual vehicle harness. Phase~A (passive) must show a physical event on the legacy bus arriving at the Clicker~4 as mirrored, timestamped T1S traffic with the Gateway never asserting TX. Phase~B (bidirectional, Mode~2 explicitly enabled) must show a Clicker-originated command reaching the legacy device and producing an observable physical response, with that response captured and mirrored back. Round-trip latency (Clicker command $\rightarrow$ observed physical response $\rightarrow$ mirrored confirmation back at Clicker) is measured per Section~\ref{sec:validation-plan}, not asserted as ``it worked.''
